% =====================================================================
%  Wstęp pracy inżynierskiej — dwie propozycje zagajenia (szkic, 26.10.2026)
%  Każdą opcję można wkleić w miejsce \section{Wstęp} / \chapter{Wstęp}.
%  \cite{TODO-...} = miejsca, gdzie trzeba podpiąć źródło z bibliografii.
%  Oryginał (opis_tematu_v2.tex) pozostaje bez zmian.
% =====================================================================

% =====================================================================
%  OPCJA A — „Skąd wiemy, z czego zbudowany jest Księżyc?”
%  (zagajenie od pytania badawczego i historii misji kosmicznych)
% =====================================================================
\chapter{Wstęp}

Skąd wiemy, z~czego zbudowany jest Księżyc? Najpełniejszą odpowiedź dały
próbki przywiezione na~Ziemię w~ramach programu Apollo -- łącznie
ok.~382~kg skał i~regolitu, których analiza trwa do~dziś~\cite{TODO-apollo}.
Sprowadzanie materiału z~powierzchni innych ciał niebieskich jest jednak
niezwykle kosztowne i~możliwe tylko w~nielicznych misjach. Dlatego
od~kilkudziesięciu lat coraz większą rolę odgrywają przyrządy, które
badają skład chemiczny skał i~gruntu bezpośrednio na~miejscu,
na~pokładzie lądownika lub łazika.

Przykładami takich przyrządów są spektrometry APXS montowane na~łazikach
marsjańskich, od~Sojournera po~Curiosity, oraz spektrometr fluorescencji
rentgenowskiej PIXL łazika Perseverance~\cite{TODO-apxs,TODO-pixl}.
Wszystkie one wykorzystują to samo zjawisko: atom pobudzony
wysokoenergetycznym promieniowaniem lub cząstką emituje promieniowanie X
o~energii właściwej tylko dla danego pierwiastka. Każdy z~nich musi
przy tym zmieścić się w~ścisłym budżecie masy, wymiarów i~poboru mocy,
który w~technice kosmicznej decyduje o~tym, czy przyrząd w~ogóle
poleci w~misję.

Na~Ziemi tę samą zasadę wykorzystuje spektroskopia EDS
(ang.~\textit{Energy Dispersive Spectroscopy}), stosowana najczęściej
w~elektronowych mikroskopach skaningowych. Wiązka elektronów pada tam
bezpośrednio na~próbkę, a~spektrometr mierzy energie emitowanych fotonów
i~na~tej podstawie ustala skład pierwiastkowy. Klasyczne działo
elektronowe z~żarzoną katodą, kolumną elektronooptyczną i~układem
próżniowym waży jednak dziesiątki kilogramów i~pobiera znaczną moc,
co wyklucza jego użycie na~łaziku.

Rozwiązaniem może być technologia MEMS
(ang.~\textit{Micro-Electro-Mechanical Systems}), pozwalająca wytwarzać
struktury elektromechaniczne o~wymiarach od~mikrometrów do~milimetrów.
W~Katedrze ... Politechniki Wrocławskiej opracowano miniaturowe źródło
promieniowania rentgenowskiego w~tej technologii, z~zimną katodą
z~nanorurek węglowych (CNT)~\cite{TODO-urbanski}. Niniejsza praca
podejmuje kolejny krok: wykorzystuje tę koncepcję do~budowy
miniaturowego analizatora EDS, w~którym wiązka elektronów trafia
bezpośrednio w~badaną próbkę. Na~powierzchni ciał pozbawionych
atmosfery, takich jak Księżyc czy planetoidy, naturalna próżnia
umożliwia propagację wiązki bez dodatkowej komory, a~źródło elektronów
o~wymiarach rzędu centymetrów i~masie rzędu gramów mogłoby stać się
częścią przyszłych przyrządów planetarnych.

Tematyka pracy wpisuje się w~dyscyplinę automatyka, elektronika,
elektrotechnika i~technologie kosmiczne, ponieważ łączy elektronikę
układów zasilania źródła elektronów i~akwizycji widm, elektrotechnikę
wysokich napięć przyspieszających i~ogniskowania wiązki oraz
projektowanie przyrządu przeznaczonego do~misji kosmicznych.


% =====================================================================
%  OPCJA B — „Laboratorium, które mieści się w dłoni”
%  (zagajenie od idei miniaturyzacji w elektronice)
% =====================================================================
\chapter{Wstęp}

Historia elektroniki jest w~dużej mierze historią miniaturyzacji.
Pierwsze komputery zajmowały całe sale, a~dziś układ o~wielokrotnie
większej mocy obliczeniowej mieści się w~telefonie. Ten sam kierunek
obrała technologia MEMS (ang.~\textit{Micro-Electro-Mechanical Systems}),
która przeniosła techniki mikroelektroniczne -- fotolitografię, trawienie
krzemu i~łączenie płytek -- na~wytwarzanie struktur
elektromechanicznych o~wymiarach od~mikrometrów do~milimetrów.
Akcelerometry, mikrofony czy czujniki ciśnienia wykonane w~tej
technologii są dziś w~każdym smartfonie~\cite{TODO-mems}.

Znacznie wolniej miniaturyzacja postępuje w~aparaturze analitycznej.
Spektroskopia EDS (ang.~\textit{Energy Dispersive Spectroscopy}),
jedna z~podstawowych metod badania składu pierwiastkowego materiałów,
wciąż kojarzy się z~elektronowym mikroskopem skaningowym ustawionym
w~laboratorium. Zasada pomiaru jest przy tym prosta: wiązka elektronów
pada na~próbkę, każdy zawarty w~niej pierwiastek emituje promieniowanie X
o~właściwej sobie energii, a~spektrometr mierzy te energie i~ustala
skład próbki. O~rozmiarach przyrządu decyduje przede wszystkim
działo elektronowe -- żarzona katoda, kolumna elektronooptyczna
i~układ próżniowy.

Zmniejszenie tego działa do~wymiarów rzędu centymetrów otwiera
zastosowania, w~których masa, wymiary i~pobór energii są kluczowe.
Najbardziej wymagającym z~nich jest technika kosmiczna. Analiza skał
i~gruntu bezpośrednio na~powierzchni Księżyca czy planetoid pozwoliłaby
poznać ich skład bez kosztownego sprowadzania próbek na~Ziemię,
a~brak atmosfery na~tych ciałach zapewnia naturalną próżnię,
w~której wiązka elektronów może się swobodnie rozchodzić.

Krokiem w~tym kierunku jest miniaturowe źródło promieniowania
rentgenowskiego w~technologii MEMS z~zimną katodą z~nanorurek węglowych
(CNT), opracowane na~Politechnice Wrocławskiej~\cite{TODO-urbanski}.
Zimna emisja polowa nie wymaga elementu grzejnego, co ogranicza pobór
mocy, a~cała struktura krzemowo-szklana waży kilka gramów. Niniejsza
praca wykorzystuje tę koncepcję do~budowy miniaturowego analizatora EDS,
w~którym wiązka elektronów, ogniskowana polem magnetycznym, trafia
bezpośrednio w~badaną próbkę.

Tematyka pracy wpisuje się w~dyscyplinę automatyka, elektronika,
elektrotechnika i~technologie kosmiczne, ponieważ łączy elektronikę
układów zasilania źródła elektronów i~akwizycji widm, elektrotechnikę
wysokich napięć przyspieszających i~ogniskowania wiązki oraz
projektowanie przyrządu przeznaczonego do~misji kosmicznych.
